\section{Equações do Movimento Orbital e Manobras Orbitais Associadas ao Problema de RVD/B: Dinâmica Orbital}
\label{sec:eq_mov_orbital_manobras}

A formulação da dinâmica orbital no contexto de missões de RVD/B leva em consideração as diferentes fases da missão.
Iniciada pelo lançamento até a execução de manobras orbitais que posicionam a espaçonave visitante em trajetória para a aproximação.
Essas fases estão relacionadas à injeção orbital, à manutenção da coplanaridade orbital do visitante com o alvo, e ao ajuste do ângulo de fase para que ocorra a transição às etapas de encontro.
Neste trabalho, a formulação matemático da dinâmica orbital  considera os vínculos de injeção inicial e as manobras de transferência, sempre de acordo com as condições ou restrições impostas pela missão de RVD/B.

\subsection{Lançamento: vínculos associados às missões de RVD/B}
\label{sec:lancamento_vinculos}

O início da missão de RVD/B ocorre na etapa de lançamento, conforme ilustrado na Figura~\ref{fig:etapas_rvdb}.
Nessa fase, a espaçonave visitante deve ser injetada em uma órbita compatível com a do alvo,para aproveitar a janela de lançamento planejada e evitar inserções em órbitas inadequadas, pois poderão demandar correções adicionais de trajetória.  

\begin{figure}[ht]
\centering
\includegraphics[width=1\linewidth]{2 Formulacao do problema/Imagens/Lançamento quadro explicativo.png}
\caption{Etapas da missão de RVD/B, adaptado de \cite{fehse2003}.}
\label{fig:etapas_rvdb}
\end{figure}

A definição da janela de lançamento deve respeitar a geometria orbital do alvo, permitindo assim, que a espaçonave visitante seja inserida em condições próximas em termos de geometria do plano orbital, e em regime estável e estabilizada.
Dessa forma, o lançamento não só coloca a espaçonave visitante em órbita, mas também estabelece as condições iniciais para a continuidade das fases subsequentes, como a órbita de faseamento, aproximação e acoplamento.